\chapter{Área Básica de Ingresso e Trajetórias Formativas}
\label{chp:abi}

Este capítulo detalha a proposta de organização do \ppccurso em uma Área Básica de Ingresso (ABI) compartilhada com o Curso Superior de Tecnologia em Automação Industrial. Trata-se de dois cursos de graduação distintos — com PPC, matriz curricular, perfil de egresso, requisitos de integralização, colação de grau e diploma próprios —, cujas matrizes possuem uma etapa formativa inicial comum. A implantação do ingresso unificado, da escolha de trajetória e de qualquer mecanismo de continuidade entre os cursos depende de ato institucional específico e de disposições simétricas nos dois PPCs.

A origem institucional desta ABI é comum aos dois cursos: ambos foram contemplados pelo Programa Universidades Inovadoras e Sustentáveis~\cite{programa_universidades_inovadoras_sustentaveis}, que concedeu à FEELT novas vagas docentes e de técnicos de laboratório — no caso deste curso, sustentando sua reformulação e renomeação a partir do curso já existente de Engenharia de Controle e Automação, com mudança para o turno noturno e ampliação de 15 para 30 vagas ofertadas (\autoref{chp:justificativa}); no caso do Curso Superior de Tecnologia, sustentando sua criação como curso inteiramente novo. Essa origem comum é o que justifica institucionalmente a organização dos dois cursos em uma mesma ABI, e não apenas a afinidade técnica entre as áreas.

O conceito que orienta esta organização curricular pode ser sintetizado da seguinte forma: o Curso Superior de Tecnologia em Automação Industrial constitui uma \textbf{formação tecnológica completa, autônoma e terminal}; o \ppccurso constitui essa mesma \textbf{formação tecnológica comum, ampliada} por fundamentação científica, matemática e metodológica própria da Engenharia, que sustenta competências de concepção, modelagem, análise, projeto, dimensionamento, integração e validação de sistemas de maior abrangência e complexidade (\autoref{sec:distincao_perfis}). A base tecnológica comum não torna os dois cursos equivalentes, nem reduz o Curso Superior de Tecnologia a uma etapa incompleta deste curso: trata-se de duas formações superiores distintas, cada uma legítima e completa segundo seus próprios objetivos, terminalidades e perfis profissionais de egresso.

\section{Caracterização da Área Básica de Ingresso}

A Área Básica de Ingresso (ABI) é uma forma de organização acadêmica em que o estudante ingressa em uma etapa inicial comum associada a mais de um curso e, posteriormente, segue uma das graduações vinculadas. A ABI não é curso autônomo, não confere titulação própria e não substitui a identidade regulatória dos cursos. As Normas Gerais da Graduação da UFU~\cite{ufu_congrad_46_2022} não regulamentam, por si só, todos os elementos operacionais desta proposta entre dois cursos distintos; por isso, a forma de ingresso, o registro acadêmico e a escolha de trajetória deverão constar de ato institucional próprio.

O compartilhamento do núcleo tecnológico inicial entre os dois cursos decorre da afinidade entre as áreas de Automação Industrial e de Controle e Automação, da racionalidade acadêmica na integração curricular entre cursos correlatos, do aproveitamento de conhecimentos comuns e do propósito institucional de constituir, para o(a) estudante, uma trajetória formativa integrada e progressiva — sem que isso implique fundir os dois cursos em um único perfil de egresso ou tratar um deles como decorrência incompleta do outro.

A escolha da trajetória ao final da etapa comum (\autoref{sec:escolha}) não deve ser interpretada como seleção entre um curso completo e um curso incompleto, tampouco como decisão entre concluir integralmente uma formação ou apenas parte dela: trata-se da escolha entre duas formações superiores diferentes, ambas legítimas, ambas completas segundo seus próprios objetivos formativos, e ambas com terminalidade profissional própria (\autoref{sec:distincao_perfis}).

A ABI não deve ser confundida com:

\begin{itemize}
    \item ciclo básico ou núcleo comum genérico, sem organização formal de ingresso e de escolha entre graduações;
    \item migração curricular — mudança de estudantes vinculados a uma matriz curricular anterior para uma nova matriz do mesmo curso (\autoref{sec:migracao_curricular});
    \item transferência de curso ou novo ingresso por processo seletivo;
    \item mudança automática de curso; e
    \item dupla diplomação automática — a obtenção de cada diploma depende do cumprimento integral dos requisitos da respectiva graduação (\autoref{sec:diplomas}).
\end{itemize}

\textbf{Pendência institucional crítica:} aprovar o ato de criação e a denominação oficial da ABI, com cursos e códigos vinculados, vagas, ingresso, registro acadêmico, procedimento de escolha, autoridade decisória e regras de transição.

\section{Distinção entre as formações e os perfis profissionais}
\label{sec:distincao_perfis}

Embora o Curso Superior de Tecnologia em Automação Industrial e o \ppccurso compartilhem um núcleo tecnológico comum, as duas formações possuem objetivos, competências, terminalidades e perfis profissionais distintos. O Curso Superior de Tecnologia em Automação Industrial constitui uma graduação tecnológica completa e autônoma, orientada à aplicação, à integração, à implementação, à operação, à supervisão, à manutenção, à gestão e ao desenvolvimento de sistemas de automação industrial. O \ppccurso incorpora essa base tecnológica e a amplia por meio de formação científica, matemática, metodológica e multidisciplinar própria da Engenharia, desenvolvendo competências para conceber, modelar, analisar, dimensionar, projetar, integrar e validar sistemas e soluções de engenharia de maior abrangência e complexidade (\autoref{chp:perfil_egresso}; \autoref{chp:objetivos}).

A ampliação promovida pelo \ppccurso não se resume a um acréscimo de carga horária, de disciplinas ou de conteúdos matemáticos: a fundamentação científica e matemática adicional sustenta o desenvolvimento de capacidades próprias de engenharia — formulação e análise de problemas complexos, modelagem matemática de sistemas físicos e dinâmicos, análise de estabilidade e de desempenho, concepção de arquiteturas de controle e automação, dimensionamento e projeto de sistemas integrados, simulação e validação de soluções, otimização de processos e recursos, análise de riscos e tomada de decisão técnica considerando aspectos econômicos, sociais, ambientais e éticos, entre outras (\autoref{chp:perfil_egresso}). O(a) engenheiro(a) formado(a) por este curso domina também os conteúdos tecnológicos desenvolvidos no Curso Superior de Tecnologia em Automação Industrial, mas os aplica em um contexto de maior profundidade científica, capacidade de abstração, integração multidisciplinar e responsabilidade técnica sobre sistemas complexos.

Essa distinção é qualitativa, não apenas quantitativa, e pode ser sintetizada como segue:

\begin{itemize}
    \item \textbf{Perfil predominante do(a) Tecnólogo(a) em Automação Industrial:} formação especializada, tecnológica e aplicada, direcionada à integração, implementação, operação, supervisão, manutenção, gestão e desenvolvimento de soluções de automação industrial em contextos delimitados e compatíveis com essa formação.
    \item \textbf{Perfil predominante do(a) engenheiro(a) formado(a) pelo \ppccurso:} formação ampla, científica, matemática, tecnológica e multidisciplinar, direcionada à análise, modelagem, concepção, dimensionamento, projeto, integração, validação e gestão de sistemas e soluções de engenharia de maior abrangência e complexidade.
\end{itemize}

Essa distinção não é absoluta: ambos os perfis podem participar de atividades de projeto, desenvolvimento e integração de sistemas de automação, porém com níveis de profundidade, abrangência, complexidade e responsabilidade técnica coerentes com a respectiva formação — não sendo adequado descrever o(a) Tecnólogo(a) como profissional que apenas executa ou opera, nem o(a) engenheiro(a) como o(a) único(a) capacitado(a) a projetar.

O quadro a seguir sintetiza essa distinção:

\begin{longtblr}[
    theme = ppc,
    caption = {Distinção entre o Curso Superior de Tecnologia em Automação Industrial e o \ppccurso},
    label = {tab:comparativo_cst_engenharia},
]{
  colspec = {Q[l,m,wd=30mm]X[l,m]X[l,m]},
  rowhead = 1,
  rowsep = 1pt,
  row{odd} = {bg=CinzaClaro},
  row{1} = {bg=AzulEscuro, fg=white},
}
\textbf{Dimensão} & \textbf{Curso Superior de Tecnologia em Automação Industrial} & \textbf{\ppccurso} \\
Natureza & Graduação tecnológica & Bacharelado em Engenharia \\
Terminalidade & Completa e autônoma & Completa e autônoma \\
Ênfase & Aplicação, integração e gestão de tecnologias de automação & Concepção, modelagem, projeto e validação de sistemas de engenharia \\
Formação científica e matemática & Aplicada às necessidades do campo tecnológico & Ampla e aprofundada, para análise e projeto de sistemas complexos \\
Complexidade predominante & Sistemas e processos tecnológicos delimitados & Sistemas multidisciplinares de maior abrangência e complexidade \\
Perfil do egresso & Tecnólogo(a) em Automação Industrial & \ppctitulacao \\
Diploma & Diploma próprio de graduação tecnológica & Diploma próprio de bacharelado em Engenharia \\
Base comum & Núcleo tecnológico integrado (\autoref{sec:etapa_comum}) & Núcleo tecnológico integrado, acrescido da formação própria de Engenharia \\
\end{longtblr}

\section{Cursos integrantes}

Integram esta Área Básica de Ingresso (ABI):

\begin{enumerate}
    \item o \ppccurso, objeto deste PPC, de grau \ppcgrau, com duração mínima de \ppctempominimo, objeto de PPC próprio; e
    \item o Curso Superior de Tecnologia em Automação Industrial, curso de graduação de grau Tecnólogo, com duração mínima de seis períodos (três anos), objeto de PPC próprio, elaborado separadamente do presente documento.
\end{enumerate}

Os dois cursos mantêm bases regulatórias distintas — o Curso Superior de Tecnologia em Automação Industrial está estruturado em conformidade com o Catálogo Nacional de Cursos Superiores de Tecnologia (CNCST, 4ª edição) e com as Diretrizes Curriculares Nacionais Gerais para a Educação Profissional e Tecnológica, ao passo que o \ppccurso se estrutura conforme as Diretrizes Curriculares Nacionais do Curso de Graduação em Engenharia~\cite{cne_ces_2_2019,cne_ces_1_2021}. A articulação pela ABI não uniformiza essas bases regulatórias: cada curso preserva seu próprio marco normativo, sua própria matriz curricular e seu próprio perfil de egresso.

\textcolor{red}{[CONFIRMAR A FORMA DE INGRESSO E O CÓDIGO ACADÊMICO INICIAL]}

\textcolor{red}{[CONFIRMAR A EXISTÊNCIA E A DISTRIBUIÇÃO DE VAGAS]}

\section{Etapa formativa comum}
\label{sec:etapa_comum}

Os cinco primeiros períodos da trajetória acadêmica na ABI são comuns entre o \ppccurso e o Curso Superior de Tecnologia em Automação Industrial. É importante destacar que esses cinco períodos \textbf{não são adicionais} à duração de nenhum dos dois cursos:

\begin{itemize}
    \item correspondem aos cinco primeiros períodos da própria matriz curricular do \ppccurso, cuja duração mínima total é de dez períodos (\ppctempominimo; \autoref{sec:regime}); e
    \item correspondem, simetricamente, aos cinco primeiros períodos da matriz curricular do Curso Superior de Tecnologia em Automação Industrial, cuja duração mínima total é de seis períodos.
\end{itemize}

Ou seja: um mesmo bloco de cinco períodos é, ao mesmo tempo, o início da trajetória do \ppccurso e o início da trajetória do Curso Superior de Tecnologia em Automação Industrial — não um acréscimo prévio a nenhuma das duas. A organização geral da ABI é, portanto:

\begin{center}
Ingresso na Área Básica de Ingresso (ABI) \\
$\downarrow$ \\
realização dos cinco primeiros períodos comuns \\
$\downarrow$ \\
escolha da graduação que o estudante pretende concluir inicialmente (\autoref{sec:escolha}) \\
$\downarrow$ \\
continuidade na trajetória específica escolhida (\autoref{sec:trajetoria_tecnologo}; \autoref{sec:trajetoria_engenharia})
\end{center}

Verificação realizada nesta revisão: os componentes curriculares ativos dos cinco primeiros períodos são atualmente idênticos, código a código, entre a matriz curricular deste curso e a do Curso Superior de Tecnologia em Automação Industrial — confirmando que a etapa comum é, de fato, integralmente compartilhada, e não apenas amplamente compartilhada.

\textcolor{red}{[A MATRIZ CURRICULAR DESTE PERFIL JÁ CADASTRA COMPONENTES CURRICULARES PRÓPRIOS DE ENGENHARIA PARA O 6º PERÍODO (405H: ESTÁGIO SUPERVISIONADO OBRIGATÓRIO I, PROJETO INTERDISCIPLINAR, PROTEÇÃO E ACIONAMENTOS INDUSTRIAIS, INSTALAÇÕES ELÉTRICAS E EXPERIMENTAL, INSTRUMENTAÇÃO E DISPOSITIVOS INTELIGENTES, E ATIVIDADES CURRICULARES DE EXTENSÃO III) E AS 16 OPTATIVAS DAS ÁREAS DE FORMAÇÃO CONTROLE E IIOT (\autoref{chp:modulos_optativos}), MAS OS PERÍODOS 7º A 10º AINDA CONTAM APENAS COM COMPONENTES ISOLADOS (ATIVIDADES CURRICULARES DE EXTENSÃO IV E V, TRABALHO DE CONCLUSÃO DE CURSO E ESTÁGIO SUPERVISIONADO OBRIGATÓRIO II — 60H, 60H, 30H E 40H, RESPECTIVAMENTE), SEM O CORPO DE DISCIPLINAS TÉCNICAS ESPECÍFICAS DE ENGENHARIA QUE COMPLETARIA ESSES PERÍODOS. PENDENTE DE ELABORAÇÃO PELO NDE/COLEGIADO ANTES DA SUBMISSÃO DESTE PPC]}

\subsection{Módulos ACE comuns e específicos}

Essa identidade curricular da etapa comum se estende às Atividades Curriculares de Extensão (ACE), organizadas em ambos os cursos em módulos curriculares de extensão (\autoref{sec:ace} de cada PPC). Os módulos \textbf{ACE I} (4º período, 90h), \textbf{ACE II} (5º período, 90h) e \textbf{ACE III} (6º período, 60h) são, código a código e hora a hora, o mesmo módulo curricular nos dois cursos, o que permite seu aproveitamento integral entre as duas matrizes em caso de permanência de vínculo (\autoref{sec:permanencia}), sem repetição de atividades extensionistas. O módulo ACE III é comum mesmo ofertado no 6º período — já o primeiro período em que as demais disciplinas das duas matrizes se distinguem (Seção~\ref{sec:etapa_comum}) —, o que faz do Curso Superior de Tecnologia em Automação Industrial um curso sem nenhum módulo ACE específico próprio: seus três módulos são integralmente compartilhados com este curso. Os módulos \textbf{ACE IV} e \textbf{ACE V} (7º e 8º período, 60h cada), sem correspondência na matriz do Curso Superior de Tecnologia, de seis períodos, são específicos deste curso, com eixos temáticos de aprofundamento em robótica, inteligência artificial e pesquisa aplicada (\autoref{sec:distincao_perfis}). O detalhamento de cada módulo — identificação, objetivos, resultados de aprendizagem, formas de integralização, avaliação, registro e governança — consta da Seção~\ref{sec:ace} de cada PPC.

\section{Escolha ao final do quinto período}
\label{sec:escolha}

Ao completar o quinto período — 2 (dois) anos e 6 (seis) meses de curso —, o estudante deverá formalizar a escolha da graduação que pretende concluir inicialmente: o \ppccurso ou o Curso Superior de Tecnologia em Automação Industrial. Essa escolha define a trajetória específica que o estudante seguirá a partir do sexto período (\autoref{sec:trajetoria_tecnologo}; \autoref{sec:trajetoria_engenharia}), mas não representa:

\begin{itemize}
    \item renúncia definitiva à outra graduação;
    \item transferência de curso;
    \item migração curricular (\autoref{sec:migracao_curricular});
    \item novo ingresso por processo seletivo;
    \item obtenção automática de dois diplomas;
    \item unificação dos dois cursos em uma única graduação; ou
    \item escolha irreversível que impeça a formação posterior na outra graduação.
\end{itemize}

O estudante concluirá a graduação escolhida. A possibilidade de prosseguir posteriormente na outra graduação somente existirá se for instituída por ato específico da UFU (\autoref{sec:permanencia}).

\textbf{Pendência institucional:} definir, antes da oferta, procedimento, prazo, órgão responsável, capacidade de cada curso e eventuais critérios acadêmicos para a escolha da graduação.

\section{Trajetória do curso superior de Tecnologia}
\label{sec:trajetoria_tecnologo}

Ao optar por concluir inicialmente o Curso Superior de Tecnologia em Automação Industrial, o estudante cursa o sexto período — específico desse curso, sem correspondência com a matriz do \ppccurso — e cumpre os demais requisitos próprios dessa formação, previstos no respectivo PPC. A trajetória do Curso Superior de Tecnologia em Automação Industrial é, portanto:

\begin{itemize}
    \item cinco períodos comuns da ABI (\autoref{sec:etapa_comum}) \textbf{mais} um período específico, totalizando seis períodos — a duração mínima já configurada para esse curso, que esta organização em ABI não altera.
\end{itemize}

O detalhamento completo dessa trajetória — componentes do sexto período, requisitos de integralização, estágio e demais exigências — consta do respectivo PPC, elaborado separadamente.

\section{Trajetória da Engenharia de Automação, Robótica e Inteligência Artificial}
\label{sec:trajetoria_engenharia}

Ao optar por concluir inicialmente o \ppccurso, o estudante passa a cursar os cinco períodos específicos desta formação (sexto ao décimo período) e a cumprir os requisitos próprios aqui previstos. A trajetória do \ppccurso é, assim:

\begin{itemize}
    \item cinco períodos comuns da ABI (\autoref{sec:etapa_comum}) \textbf{mais} cinco períodos específicos, totalizando dez períodos — a duração mínima de \ppctempominimo.
\end{itemize}

O Estágio Supervisionado Obrigatório (dividido em duas etapas, no 6º e no 10º período), o Trabalho de Conclusão de Curso (9º período) e o Projeto Interdisciplinar (6º período) já constam da matriz curricular deste perfil e são detalhados na Seção~\ref{sec:acc_conclusao} do \autoref{chp:estrutura_curricular}. \textcolor{red}{[DETALHAR AS DEMAIS DISCIPLINAS TÉCNICAS ESPECÍFICAS DE ENGENHARIA DO 7º AO 10º PERÍODO — AINDA NÃO CADASTRADAS NA MATRIZ CURRICULAR DESTE PERFIL ALÉM DOS COMPONENTES JÁ MENCIONADOS E DAS 16 OPTATIVAS DAS ÁREAS DE FORMAÇÃO (\autoref{chp:modulos_optativos})]}

\section{Conclusão da primeira graduação}
\label{sec:conclusao_primeira}

Qualquer que seja a graduação escolhida ao final do quinto período (\autoref{sec:escolha}), sua conclusão pressupõe o cumprimento integral dos requisitos previstos no respectivo PPC — não há dispensa ou abreviação de requisitos em razão da organização em ABI. Para o \ppccurso, esses requisitos incluem, além dos componentes curriculares obrigatórios e optativos, os demais elementos exigidos pelas Diretrizes Curriculares Nacionais do Curso de Graduação em Engenharia~\cite{cne_ces_2_2019,cne_ces_1_2021} — estágio curricular supervisionado, trabalho de conclusão de curso e atividades complementares, entre outros.

Os requisitos de integralização próprios deste curso — Estágio Supervisionado Obrigatório I e II, Trabalho de Conclusão de Curso, Projeto Interdisciplinar e Atividades Acadêmicas Complementares — já constam da matriz curricular deste perfil (\autoref{chp:estrutura_curricular}). \textcolor{red}{[CONFIRMAR AS DEMAIS EXIGÊNCIAS DAS DIRETRIZES CURRICULARES NACIONAIS DE ENGENHARIA AINDA NÃO DETALHADAS NESTE PERFIL, EM PARTICULAR O CORPO DE DISCIPLINAS TÉCNICAS ESPECÍFICAS DO 7º AO 10º PERÍODO]}

\section{Permanência de vínculo}
\label{sec:permanencia}

As Normas Gerais da Graduação da UFU~\cite{ufu_congrad_46_2022} preveem permanência de vínculo para conclusão de outra habilitação, modalidade ou certificado de estudos oferecido pelo mesmo curso. Como esta proposta articula dois cursos com PPCs, códigos e diplomas distintos, essa norma não basta para autorizar a continuidade aqui descrita. Até a aprovação de ato institucional específico, não se assegura permanência de vínculo, ingresso na segunda graduação sem novo processo seletivo ou segundo diploma.

A eventual regulamentação deverá estabelecer que o mecanismo:

\begin{itemize}
    \item é uma possibilidade acadêmica, não uma etapa obrigatória da trajetória do estudante;
    \item constitui procedimento formal, mediante solicitação do próprio estudante — não é concedida de ofício;
    \item deve ser solicitada antes da colação de grau da primeira graduação concluída;
    \item está sujeita às normas acadêmicas institucionais vigentes, podendo depender da existência de vaga e de análise pelas instâncias competentes;
    \item permite que a continuidade na outra graduação ocorra sem caracterizar novo ingresso, migração curricular (\autoref{sec:migracao_curricular}) ou transferência de curso, observadas as normas institucionais aplicáveis a essa articulação entre dois cursos distintos; e
    \item não produz, por si só, o diploma da segunda graduação (\autoref{sec:continuidade_segunda}) — este depende do cumprimento integral dos requisitos curriculares ainda pendentes.
\end{itemize}

\textbf{Pendência institucional crítica:} definir a base normativa, a viabilidade administrativa, as vagas, o prazo, a forma do pedido, o órgão competente, os requisitos de deferimento, o registro acadêmico e os efeitos sobre colação de grau e diploma entre dois cursos distintos.

\section{Continuidade na segunda graduação}
\label{sec:continuidade_segunda}

Somente se houver regulamentação institucional e deferimento nos termos do futuro ato, o estudante poderá dar continuidade à outra graduação. Nessa hipótese:

\begin{itemize}
    \item os componentes curriculares comuns já integralizados são aproveitados, sem necessidade de nova integralização;
    \item aplicam-se as equivalências curriculares previamente definidas entre as duas matrizes (\autoref{sec:equivalencias_abi});
    \item o estudante deverá cursar os componentes específicos da segunda graduação ainda não integralizados e cumprir todos os demais requisitos curriculares dessa formação;
    \item o tempo necessário para a conclusão depende dos componentes e requisitos ainda pendentes, e não é automaticamente igual à duração da trajetória específica descrita nas Seções~\ref{sec:trajetoria_tecnologo} e~\ref{sec:trajetoria_engenharia} para quem a escolhe como primeira formação; e
    \item a conclusão da segunda graduação é seguida de nova colação de grau e da expedição do respectivo diploma (\autoref{sec:diplomas}).
\end{itemize}

A segunda graduação não é apresentada, em nenhuma hipótese, como diploma automático, habilitação adicional automática, apostilamento do primeiro diploma ou dupla diplomação automática — depende, em todos os casos, do cumprimento integral dos requisitos específicos ainda pendentes da respectiva formação.

\subsection{Do Curso Superior de Tecnologia em Automação Industrial para a Engenharia de Automação, Robótica e Inteligência Artificial}

Na hipótese condicionada ao ato institucional, o estudante que concluir inicialmente o Curso Superior de Tecnologia em Automação Industrial e obtiver autorização para a continuidade poderá prosseguir no \ppccurso, sujeito ao aproveitamento formal de estudos e aos requisitos da segunda formação. O fluxo pretendido é:

\begin{center}
Ingresso na ABI $\to$ cinco períodos comuns $\to$ escolha pelo Curso Superior de Tecnologia $\to$ sexto período específico desse curso $\to$ integralização do Curso Superior de Tecnologia $\to$ permanência de vínculo (antes da colação de grau) $\to$ colação de grau no Curso Superior de Tecnologia $\to$ continuidade no \ppccurso $\to$ aproveitamento dos componentes comuns e das equivalências $\to$ componentes específicos restantes deste curso $\to$ integralização deste curso $\to$ segunda colação de grau $\to$ diploma de \ppcgrau em Engenharia de Automação, Robótica e Inteligência Artificial.
\end{center}

\textcolor{red}{[CONFIRMAR A TABELA DE EQUIVALÊNCIAS ENTRE OS DOIS CURSOS]}

\subsection{Da Engenharia de Automação, Robótica e Inteligência Artificial para o Curso Superior de Tecnologia em Automação Industrial}

Na hipótese condicionada ao ato institucional, o estudante que concluir inicialmente o \ppccurso e obtiver autorização para a continuidade poderá prosseguir no Curso Superior de Tecnologia em Automação Industrial, sujeito ao aproveitamento formal de estudos. O detalhamento dos componentes e requisitos residuais consta do PPC desse curso.

\section{Equivalências e aproveitamento de estudos}
\label{sec:equivalencias_abi}

O reconhecimento de componentes curriculares entre o \ppccurso e o Curso Superior de Tecnologia em Automação Industrial, para os fins descritos neste capítulo, apoia-se em dois conceitos distintos:

\begin{description}
    \item[Equivalência curricular] Reconhecimento formal de compatibilidade entre componentes curriculares das duas matrizes, considerando conteúdo, carga horária, competências e resultados de aprendizagem, entre outros critérios institucionais.
    \item[Aproveitamento de estudos] Procedimento acadêmico pelo qual componentes curriculares já integralizados são reconhecidos com base nas equivalências previamente aprovadas e nas normas institucionais vigentes, dispensando o estudante de cursá-los novamente.
\end{description}

Como os cinco primeiros períodos são comuns às duas matrizes (\autoref{sec:etapa_comum}), o aproveitamento dos componentes dessa etapa, ao se solicitar a permanência de vínculo, decorre diretamente dessa identidade curricular. O aproveitamento não constitui dispensa genérica ou automática de todos os requisitos da segunda graduação: os componentes específicos ainda não integralizados, e os demais requisitos próprios da formação, permanecem exigíveis (\autoref{sec:continuidade_segunda}).

\textcolor{red}{[CONFIRMAR A TABELA DE EQUIVALÊNCIAS ENTRE OS DOIS CURSOS]}

\section{Colações de grau e expedição dos diplomas}
\label{sec:diplomas}

A conclusão de cada graduação integrante da ABI é seguida de colação de grau e de expedição de diploma próprios e distintos:

\begin{itemize}
    \item a conclusão do \ppccurso confere o diploma de \ppcgrau em Engenharia de Automação, Robótica e Inteligência Artificial, mediante colação de grau específica deste curso;
    \item a conclusão do Curso Superior de Tecnologia em Automação Industrial confere o diploma de Tecnólogo nessa formação, mediante colação de grau específica desse curso, nos termos do respectivo PPC.
\end{itemize}

Na hipótese de a continuidade ser regulamentada e de o estudante concluir as duas graduações, haverá colação de grau e diploma próprios para cada curso, conforme o futuro ato institucional e os procedimentos de registro acadêmico. Não há emissão automática ou simultânea dos dois diplomas, nem diploma único que os represente conjuntamente.

\textcolor{red}{[CONFIRMAR PROCEDIMENTO PARA A SEGUNDA COLAÇÃO DE GRAU]}

\textcolor{red}{[CONFIRMAR O REGISTRO ACADÊMICO DA SEGUNDA GRADUAÇÃO]}

\section{Implantação e inexistência de migração curricular}
\label{sec:migracao_curricular}

O \ppccurso é curso novo — encontra-se, no momento da elaboração deste PPC, em fase de proposta —, e o Curso Superior de Tecnologia em Automação Industrial, igualmente, é proposta de curso novo. A organização em ABI descrita neste capítulo é, portanto, a forma de ingresso e de articulação entre os dois cursos desde sua implantação, e não uma migração de estudantes de uma matriz curricular anterior para a estrutura aqui descrita — distinção relevante entre dois conceitos que não devem ser confundidos:

\begin{description}
    \item[Migração curricular] Mudança de estudantes vinculados a um currículo anterior de um mesmo curso para um currículo novo desse curso — situação que não se aplica a nenhum dos dois cursos integrantes desta ABI, por serem ambos cursos novos, sem currículo anterior ao aqui descrito.
    \item[Continuidade proposta entre cursos] Mecanismo ainda dependente de ato institucional específico, distinto da migração curricular e da permanência de vínculo já disciplinada para habilitações, modalidades ou certificados de um mesmo curso (\autoref{sec:permanencia}).
\end{description}

Não há migração ou transferência automática entre os cursos. A etapa curricular comum não cria, por si só, direito de ingresso na outra graduação; qualquer continuidade dependerá da regulamentação indicada na \autoref{sec:permanencia}.
